\answer
\begin{enumerate}

%%--- 章末問題 3.1 ---
\item[\ref{ex3.1}]
(a) シリコン（$\mathcal{E}_c = 3\times10^7$\,V/m）：式\ref{eq3.4}より
\begin{equation*}
W = \frac{2\times3300}{3\times10^7} = 2.2\times10^{-4}\,\mathrm{m} = 220\,\mathrm{\mu m}
\end{equation*}
式\ref{eq3.7}より
\begin{equation*}
R_{\mathrm{on}}A = \frac{4\times3300^2}{1.04\times10^{-10}\times0.135\times(3\times10^7)^3}
\approx 1.1\times10^{-4}\,\Omega\cdot\mathrm{m^2} = 1.1\,\Omega\cdot\mathrm{cm^2}
\end{equation*}
4H-SiC（$\mathcal{E}_c = 2.5\times10^8$\,V/m，$\varepsilon = 8.6\times10^{-11}$\,F/m，
$\mu_n = 0.10\,\mathrm{m^2/(V\cdot s)}$）：
\begin{equation*}
W = \frac{2\times3300}{2.5\times10^8} \approx 26\,\mathrm{\mu m}
\end{equation*}
\begin{equation*}
R_{\mathrm{on}}A = \frac{4\times3300^2}{8.6\times10^{-11}\times0.10\times(2.5\times10^8)^3}
\approx 3.2\,\mathrm{m\Omega\cdot cm^2}
\end{equation*}

(b) シリコンでは$1\,\Omega\cdot\mathrm{cm^2}$を超え，
現実的なチップ面積では大電流を流せるオン抵抗にならない。
そこでシリコン時代は，伝導度変調でドリフト層の抵抗を桁違いに下げられる
IGBT（バイポーラデバイス）が3.3\,kV級を独占してきた。
代償はテール電流によるスイッチングの遅さである。
SiCなら同じ耐圧をユニポーラのMOSFETで数m$\Omega\cdot$cm$^2$級で実現できるため，
テール電流のない高速スイッチングと低導通損失を両立でき，
鉄道用変換器などの小型化・高効率化が進んでいる。

%%--- 章末問題 3.2 ---
\item[\ref{ex3.2}]
(a) 導通損失は
$P_{\mathrm{cond}} = V_{\mathrm{CE(sat)}}ID = 1.5\times100\times0.5 = 75$\,W。
スイッチング損失は式\ref{eq3.10}より
\begin{equation*}
P_{\mathrm{sw}} = \frac{1}{2}\times400\times100\times1\times10^{-6}\times10^4 = 200\,\mathrm{W}
\end{equation*}
総損失は$P_{\mathrm{total}} = 75 + 200 = 275$\,W。

(b) $P_{\mathrm{sw}}$は周波数に比例するから$200\times5 = 1000$\,Wとなり，
$P_{\mathrm{total}} = 1075$\,W。
損失は4倍近くに跳ね上がり，冷却が現実的でなくなる。
IGBTのインバータが数十kHz以下で運転される理由である。

(c) 導通損失は
$P_{\mathrm{cond}} = R_{\mathrm{on}}I^2D = 0.01\times100^2\times0.5 = 50$\,W。
スイッチング損失は
\begin{equation*}
P_{\mathrm{sw}} = \frac{1}{2}\times400\times100\times100\times10^{-9}\times5\times10^4
= 100\,\mathrm{W}
\end{equation*}
総損失は$150$\,Wで，(b)の約7分の1である。
切り替え時間が10分の1になったことで，
周波数を5倍に上げてなお損失を減らせる。
これがワイドギャップ半導体による高周波化の定量的な意味である。

%%--- 章末問題 3.3 ---
\item[\ref{ex3.3}]
(a)
\begin{equation*}
P = \frac{1}{2}C_{\mathrm{oss}}V^2f_{\mathrm{sw}}
= \frac{1}{2}\times100\times10^{-12}\times800^2\times10^5 = 3.2\,\mathrm{W}
\end{equation*}

(b)
\begin{equation*}
P_g = Q_gV_gf_{\mathrm{sw}} = 100\times10^{-9}\times15\times10^5 = 0.15\,\mathrm{W}
\end{equation*}

(c) 例題\ref{rei3.2}の電圧・電流の重なりによる損失（$P_{\mathrm{sw}} = 5$\,W）が
最も大きく，次いで出力容量の損失（3.2\,W），
ゲート駆動損失（0.15\,W）の順である。
ただし出力容量の損失は$V^2$に比例するため，
高電圧・高周波・軽負荷（$I$が小さい）の条件では支配的になりうる。
ゲート駆動損失は素子自体ではなく駆動回路で消費されるが，
これも周波数に比例して増える点は同じである。

%%--- 章末問題 3.4 ---
\item[\ref{ex3.4}]
\rev{\ref{sec3.4}節の手順（1.耐圧 → 2.電流 → 3.損失 → 4.駆動 → 5.熱）の順にたどる。}

(a) \textbf{MOSFET（GaN HEMTを含む）}。
\rev{1.~耐圧：整流後の直流は約$141$\,V（入力変動を見て$200$\,V弱）だが，
ACアダプタは6章の絶縁型回路（フライバック）で，スイッチにはトランスの跳ね返り電圧が上乗せされる。
$600$〜$650$\,V級の定格なら50〜80\,\%の範囲に収まる。
2.~電流：$65$\,Wを$141$\,Vで割った平均電流は$0.5$\,A程度，ピークでも数Aで，
$I^* = V_{\mathrm{CE(sat)}}/R_{\mathrm{on}}$（十A級）よりはるかに小さい ― 抵抗性のMOSFETの領域である（\ref{fig3.5}）。
3.~損失：数百kHzでは切替え時間$1\,\mathrm{\mu s}$級のIGBTのスイッチング損失（式\ref{eq3.10}）が出力の半分に達し，使えない。
切替え時間数十nsのMOSFETなら数W，GaN HEMTならさらに小さい。
4.~駆動：電圧駆動で，ゲート電荷も小さく（GaNでは数nC），駆動回路は小さく損失も無視できる。
5.~熱：損失が$2$\,W程度なら，放熱器なしで基板の銅箔から逃がす熱抵抗（$30$\,K/W程度）でも
式\ref{eq3.14}の温度上昇は$60$\,Kで，筐体内が$50\,^{\circ}$Cでも$T_j \approx 110\,^{\circ}$Cと定格内に収まる。
小型化を重視するならGaN HEMTが最有力である。}

(b) \textbf{SiC-MOSFETまたはIGBT}。
\rev{1.~耐圧：電池$800$\,Vにサージを見込み，$1200$\,V級（$800/1200 \approx 67\,\%$）を選ぶ。
2.~電流：$150$\,kWを$800$\,Vで割ると$200$\,A近く，相電流は数百Aに達し，$I^*$を大きく超える ―
バイポーラ系（IGBT）の領域で，MOSFETを使うならチップを並列にする。
3.~損失：IGBTは$10$〜$20$\,kHzが限度で，章末問題\ref{ex3.2}の条件では1素子$275$\,W。
SiC-MOSFETなら同じ耐圧でテール電流がなく，$50$\,kHzでも$150$\,Wと少ない。
4.~駆動：どちらも電圧駆動で，ゲート電荷が数百nCあっても駆動損失は$1$\,W以下にすぎない。
5.~熱：損失が百W級なので自然空冷では捨てられず，\ref{sec3.4}節の例のように
水冷で熱抵抗を$0.5$\,K/W程度まで下げて$T_j$を定格内に収める。
損失の少ないSiC-MOSFETは，冷却器と電池を小さくできる分だけ車両全体で有利になる。}

(c) \textbf{サイリスタ}。
\rev{1.~耐圧：数百kVは1素子の耐圧（$10$\,kV級が最大）をはるかに超えるので，
多数を直列にしたバルブとして使う。直列にできる素子であることが第一条件になる。
2.~電流：GW級を数百kVで割ってもkA級で，$I^*$を桁で超える ― 最大の電流容量をもつサイリスタの領域である。
3.~損失：商用周波数（$50$/$60$\,Hz）ではスイッチング損失は無視でき，損失はほぼ導通損失だけになる。
4.~駆動：電流駆動で，ゲートではオフできないが，商用周波数の交流に同期して電流が自然にゼロになるので他励形で運用できる。
高電位の素子へは光でゲート信号を送る。
5.~熱：1素子あたりの導通損失はkW級に達するので，純水で冷却して熱抵抗を数十分の1\,K/Wまで下げる。
電流駆動と転流の制約は，この電力規模では受け入れるしかない
（近年は，IGBTを多段に重ねてゲートでオフできるようにした自励式の変換所も増えている）。}

%%--- 章末問題 3.5 ---
\item[\ref{ex3.5}]
シリコンを基準にすると，バリガ性能指数の比は
\begin{equation*}
\frac{\mathrm{BFOM_{SiC}}}{\mathrm{BFOM_{Si}}}
= \frac{9.7}{11.7}\times\frac{1000}{1350}\times\left(\frac{2.5\times10^8}{3\times10^7}\right)^3
\approx 0.83\times0.74\times579 \approx 360
\end{equation*}
\begin{equation*}
\frac{\mathrm{BFOM_{GaN}}}{\mathrm{BFOM_{Si}}}
= \frac{9.0}{11.7}\times\frac{1200}{1350}\times\left(\frac{3.3\times10^8}{3\times10^7}\right)^3
\approx 0.77\times0.89\times1331 \approx 910
\end{equation*}
理論上，同じ耐圧の特性オン抵抗はSiCで約360分の1，GaNで約900分の1になる。

市販デバイスがこの比率ほど下がっていないのは，
式\ref{eq3.7}がドリフト層だけの抵抗だからである。
実デバイスにはMOS界面のチャネル抵抗（SiCでは界面の欠陥により移動度が
バルク値より大幅に低い），基板やパッケージの抵抗が直列に加わる。
また信頼性の余裕を見て理論限界より厚め・低濃度めに設計される。
それでも限界線が2桁半下にあるという事実が，
改良の余地がまだ大きいことを保証している。

\end{enumerate}
